\section{Completion of the proof}\label{sec:conclusion}

We now use Proposition~\ref{prop:comparisons} to control the diagonals in
the transfers and then symmetrize the final amplitude.  All parameters,
priors, and zero conventions are those of Section~\ref{sec:construction}.
In particular, the constants denoted by $C$ in costs $Crm$ are independent
of the depth $k$ and of the large gap constants.  The limits are taken with
all these constants and $k$ fixed, and then $L\to\infty$.

\subsection{Counting bad arrangements}

At level $l$ put $r=2^l$.  Call a pair of arrangements of the $rm$ bulk
variables \emph{bad} if it is not covered by \eqref{eq:src56}.  Thus every
pair is called bad when $l=0,1$; when $l\ge2$, a bad pair has more than
$3r/4$ connected components in its bipartite overlap graph.  For each
fixed first arrangement, the number of bad second arrangements satisfies
\begin{equation}\label{eq:src60}
 N_{\mathrm{bad}}(r,m)
 \le (m!)^r
       \exp\bigl((\tfrac14\log r+C)rm+O_k(1)\bigr).
\end{equation}
Here arrangements distinguish all bulk positions and all sampled bulk
variables, so they are indexed by permutations even before any
distinctness restrictions are imposed.

To prove \eqref{eq:src60} for $l\ge2$, write $a_i$ for the number of leaves
on either side of the $i$th connected component.  The numbers on the two
sides agree: every vertex has degree $m$, and counting the component's
edges on either side gives the equality.  If there are $t$ components,
then
\[
 \sum_{i=1}^t a_i=r,
 \qquad
 \sum_{i=1}^t(a_i-1)=r-t<\frac r4.
\]
There are only $\exp(O_k(1))$ ways to choose and pair the two partitions
of the leaf sets into these components.  Once they are chosen, the
variables incident to a component can be assigned to its second set of
slots in at most $(a_i m)!$ ways.  For $1\le a\le r$,
\[
 a\log a
  =(a-1)\log a+\log a
  \le(a-1)(\log r+1),
 \qquad
 (am)!\le(m!)^a a^{am}.
\]
The factorial inequality follows by bounding one multinomial
coefficient by the sum of all multinomial coefficients with $a$ parts.
Consequently
\[
 \prod_i(a_i m)!
 \le(m!)^r\exp\left(m\sum_i a_i\log a_i\right)
 \le(m!)^r\exp\left(\frac14(\log r+1)rm\right),
\]
which proves the claim.  For $r=1,2$, the upper bound
$(rm)!\le(m!)^r r^{rm}$ proves the same assertion after enlarging the
absolute constant $C$.  Nonbulk matchings in a diagonal have only
$\exp(O_k(1))$ possibilities, since their number of slots is fixed once
$k$ is fixed.

We shall also use the corresponding bound for the fraction of bad
arrangements.  Since $m!\le m^m$ and $(rm)!\ge(rm/e)^{rm}$,
\begin{equation}\label{eq:conclusion-bad-fraction}
 \frac{N_{\mathrm{bad}}(r,m)}{(rm)!}
 \le\exp\bigl((-\tfrac34\log r+C)rm+O_k(1)\bigr).
\end{equation}

\subsection{Diagonal bounds and the order of parameter choices}

Let $\mathcal D_j$ be the nonnegative diagonal contribution in the
extended square at step $j$, before the factor $e^{c_g}$ in
\eqref{eq:src53}.  At this step the two compared amplitudes have level
$j-1$ and $r_j=2^{j-1}$ leaves.  We first estimate the bad matchings.

For each such matching, apply
$|A^{(1)}\overline{A^{(2)}}|\le
(|A^{(1)}|^2+|A^{(2)}|^2)/2$ on their simultaneous support, after the
nonnegative-square reduction that gives \eqref{eq:src54}.  For either
retained square, extract the point probability of the \emph{other}
ordered tuple.  On this support it is at most
\[
 C_H\exp(-T_H+O_k(1)),
 \qquad
 C_H\le L^{-r_jm}\exp(Cr_jm+O_k(1)).
\]
The retained tuple is then summed with its own priors.  This order of
extraction is valid for both squares, including when a matching exchanges
positions with different log-cell priors.  The external integer measure
and the $u$ weight contribute $\exp(G+r_jw_j)$; by
\eqref{eq:src48}, their combination with $\exp(-T_H)$ is
$\exp(-\Delta_j)$.  All remaining Archimedean factors are bounded by
$\exp(O_k(1))$ on the simultaneous support.  Thus the one-assignment
estimate \eqref{eq:src55}, followed by \eqref{eq:src60}, bounds the bad
part by
\[
 \exp(-\Delta_j) L^{-r_jm}(m!)^{r_j}
 \exp\bigl(r_j\Delta_0+
       (\tfrac14\log r_j+C)r_jm+o_k(m)\bigr).
\]
Since $m/L\le z$, we have
$L^{-r_jm}(m!)^{r_j}\le\exp(r_jm\log z)$.  We obtain
\begin{equation}\label{eq:src61}
 \mathcal D_{j,\mathrm{bad}}
 \le \exp\left(
 -\Delta_j+r_j\bigl\{\Delta_0+
       (\log z+\tfrac14\log r_j+C)m\bigr\}+o_k(m)
 \right).
\end{equation}
For every other matching, \eqref{eq:src56} applies with its stated
Archimedean multiplier.  After the probability extraction, the number
and size of these terms cost at most $\exp(O_k(m))$: in particular,
\[
 (r_jm)!L^{-r_jm}
 \le\exp\bigl(r_jm\log(r_jm/L)\bigr)
 =\exp(O_k(m)).
\]
Their total is therefore $\exp(-\omega_k(m))$.  In this notation,
$\omega_k(m)/m\to\infty$ at every fixed $k$.  Combining the two parts,
\begin{equation}\label{eq:conclusion-diagonal-total}
 \mathcal D_j\le
 \exp\left(
 -\Delta_j+r_j\bigl\{\Delta_0+
       (\log z+\tfrac14\log r_j+C)m\bigr\}+o_k(m)
 \right)+\exp(-\omega_k(m)).
\end{equation}

We specify the parameter order and the reserve in the iteration.  The
constant $B_s$ has already been chosen sufficiently large for
\eqref{eq:src47}, whose constant $B_*$ is independent of the later gaps
and of $k$.  Write
\[
 Z_G=\sum_{p\ \mathrm{prime}}
             \frac{\varphi(\log p-G)}p,
 \qquad c_g=\log Z_G^{-1}.
\]
The giant range and the prime-cell normalization give
$Z_G\sim G^{-1}$ and $c_g\le .91L$ for all sufficiently large $L$,
with the coefficient $.91$ independent of all the gap choices and $k$.
Fix, now,
\begin{equation}\label{eq:conclusion-budget}
 B=B_*+3.
\end{equation}
Choose $B_z\ge9$ and then choose $B_D$ so large that
\begin{equation}\label{eq:conclusion-gap-choice}
 B_D\ge B_s+2B+C+6,
\end{equation}
where $C$ dominates the uniform constant in
\eqref{eq:conclusion-diagonal-total}.  These are fixed constants, chosen
before $k$.

Indeed, substitution from \eqref{eq:src43} shows that the exponent of
the first term in \eqref{eq:conclusion-diagonal-total}, divided by
$r_jm$, is
\begin{equation}\label{eq:conclusion-diagonal-rate}
 B_s-B_D+(9-B_z)\log z
       -\frac3{20}\log r_j+C+o_k(1).
\end{equation}
Hence, for every sufficiently large fixed $k$ and then sufficiently
large $L$, uniformly for $1\le j\le k$,
\begin{equation}\label{eq:conclusion-diagonal-reserve}
 \mathcal D_j\le\exp\bigl(-(2B+3)r_jm\bigr).
\end{equation}
The precise choice of $k$ will be made below.  We require already that
it be large enough that, using $m=k^4L+O(1)$,
\begin{equation}\label{eq:conclusion-cg-reserve}
 \frac{c_g+\log2}{m}\le\frac2{k^4}<1
\end{equation}
for all sufficiently large $L$.

The transfer identity underlying \eqref{eq:src53} gives
\[
 e^{-c_g}|\eta_{j-1}|^2
 \le\mathcal D_j+\eta_j
 \le\mathcal D_j+|\eta_j|.
\]
If $|\eta_{j-1}|\ge\exp(-Br_jm)$, then
\eqref{eq:conclusion-diagonal-reserve} and
\eqref{eq:conclusion-cg-reserve} imply
$\mathcal D_j\le\tfrac12 e^{-c_g}|\eta_{j-1}|^2$.  Therefore
\begin{equation}\label{eq:conclusion-transfer-lower}
 |\eta_j|\ge\frac12e^{-c_g}|\eta_{j-1}|^2.
\end{equation}
Starting with \eqref{eq:src47}, induction proves simultaneously this
estimate and the sharper bound
\begin{equation}\label{eq:conclusion-exact-recurrence}
 -\log|\eta_j|
 \le 2^jB_*m+(2^j-1)(c_g+\log2)
 <B2^jm,
 \qquad 0\le j\le k.
\end{equation}
For the induction, the sharper bound at $j-1$ supplies the coarse
hypothesis needed for \eqref{eq:conclusion-transfer-lower}; that
inequality then gives the displayed sharper bound at $j$.  Its final
inequality follows from \eqref{eq:conclusion-budget} and
\eqref{eq:conclusion-cg-reserve}.  In particular the reserve is inherited
from the initial estimate rather than spent afresh at each step.

The transfers have now retained an exponential lower bound for
$|\eta_k|$.  We will obtain the opposite bound by averaging $A_k$ over
the bulk arrangements.  Since $G_R(s/d)$ depends only on their prime
values, not their positions, it is common to these arrangements.
The Cauchy--Schwarz step therefore also requires a norm estimate for
this regular-transform factor.

\subsection{The norm of the common regular transform}

At the final level put $r=2^k$.  The sum over $s$ in this subsection is
over the nonzero signed integers with $|s|\le V_k$.  We claim
\begin{equation}\label{eq:conclusion-regular-norm}
 \E_{\mathbf d,\mathcal I_k}
       \sum_{0<|s|\le V_k}|G_R(s/d)|^2\ll V_k,
\end{equation}
with an absolute implied constant for sufficiently large $L$ at the
fixed parameters.  Only the pairwise-distinctness and coprimality zero
conventions of $G_R$ are needed here; bins and all history restrictions
belong to the amplitudes.

Fix $s$, the spectator list $\mathbf d$, the actual regular small
primes, and one of the two giant primes, denoted by $q$.  Configurations
with a repeated regular small prime or with a regular small prime
dividing $d$ already give zero and can be discarded.  Let $C$ be the
squarefree product of the remaining regular small primes, and call the
other giant $p$.  Both giant transforms have absolute value at most
one.  After dropping their squared absolute values, the remaining
factor is
\begin{equation}\label{eq:conclusion-residue-test}
 F(p\bmod C),\qquad
 F(a)=\prod_{\ell\mid C}
       \left|g_\ell(a_\ell a^{-1}\bmod\ell)\right|^2,
 \qquad
 a_\ell=\frac{s}{dq(C/\ell)}\pmod\ell .
\end{equation}
All $a_\ell$ are units: the retained primes are distinct and coprime to
$d$, and $0<|s|\le V_k$ is smaller than every actual slot prime.  The
restriction $p\ne q$ may now be dropped, since the remaining
nonnegative expression is still defined when $p=q$.  The giant range
is disjoint from all the small-prime ranges, so every prime $p$ in its
cell is a unit modulo $C$.

Write $U_C=(\Z/C\Z)^\times$.  By the Chinese remainder theorem, as $a$
varies uniformly over $U_C$, the arguments
$a_\ell a^{-1}\pmod\ell$ are independent uniform nonzero residues.
The exact local normalizations $g_\ell(0)=0$ and
$\ell^{-1}\sum_x|g_\ell(x)|^2=1$ therefore give the identities
\begin{equation}\label{eq:conclusion-crt-norm}
 \sum_{a\in U_C}F(a)=C,
 \qquad
 \frac1{\phi(C)}\sum_{a\in U_C}F(a)
   =\frac C{\phi(C)}
   =\prod_{\ell\mid C}\frac\ell{\ell-1}.
\end{equation}
Here $\phi$ is Euler's totient.  There are $O_k(m)$ factors, and every
such prime is at least $\exp(\exp(.004L))$.  Consequently
$C/\phi(C)=1+o_k(1)$ uniformly in the fixed lists.  This calculation
uses no uniform pointwise bound on the small-prime transforms.

We next justify the averaging over the actual prime $p$, including the
possible exceptional term.  The size bounds for the final list give
\[
 \log C=O_k(h+m\exp(.006L))\le\exp(.012L)
\]
for sufficiently large $L$, so the giant row of \eqref{eq:src57}
applies with modulus $C$.  With the positive constant $c$ from that
estimate, put
\[
 \varepsilon_L=\exp(-c\exp(.018L)).
\]
Splitting the support of $\varphi(t-G)$ into
intervals of length at most one and applying partial summation gives,
uniformly for $a\in U_C$,
\begin{align}
 \sum_{\substack{p\ \mathrm{prime}\\p\equiv a\pmod C}}
       \frac{\varphi(\log p-G)}p
 &=\frac1{\phi(C)}\int
       \frac{\varphi(t-G)}t
       \bigl(1-\chi_*(a)e^{(\beta_*-1)t}\bigr)\,dt
       +O(\varepsilon_L).
       \label{eq:conclusion-progression-average}
\end{align}
As in \eqref{eq:src57}, the exceptional term is omitted if it is
absent or its conductor does not divide $C$.  When it is present,
\[
 0\le1-\chi_*(a)e^{(\beta_*-1)t}\le2.
\]
This pointwise inequality suffices even if the exceptional character
correlates with $F$.

Multiply \eqref{eq:conclusion-progression-average} by $F(a)$, sum over
$a\in U_C$, and divide by $Z_G$.  Since $F\ge0$, the exact sum in
\eqref{eq:conclusion-crt-norm} both bounds the main term and sums the
absolute progression errors.  It follows that
\begin{equation}\label{eq:conclusion-prime-norm}
 \E_p F(p\bmod C)
 \le
 2\frac C{\phi(C)}
       \frac{\displaystyle\int\varphi(t-G)\,dt/t}{Z_G}
       +O(Z_G^{-1}C\varepsilon_L)
 =2+o_k(1).
\end{equation}
Indeed the integral divided by $Z_G$ is $1+o(1)$ by ordinary prime-cell
normalization.  Moreover $Z_G^{-1}=\exp(O(L))$ and
$\log C\le\exp(.012L)$, whereas
$\varepsilon_L=\exp(-c\exp(.018L))$, so the summed error is $o(1)$.
All estimates are uniform in $s$, $\mathbf d$, $q$, and the retained
small-prime list.  Averaging those variables and summing over the at
most $2V_k$ frequencies proves \eqref{eq:conclusion-regular-norm}.

\subsection{Symmetrization and contradiction}

Let $\mathfrak S_{rm}$ permute the values in all $rm$ bulk positions,
keeping every other position fixed.  The bulk priors are identical,
including their allowed prime deletions, so their joint law is
invariant under this action.  Also $G_R(s/d)$ is unchanged: the product
$R$, the set of its prime factors, the exact transform attached to each
bulk prime value, and every argument $s/(d(R/\ell))$ are unchanged.
Its distinctness and coprimality zero conventions are invariant as
well.  Define
\begin{equation}\label{eq:conclusion-symmetrization}
 A_{\mathrm{sym}}(s)
   =\frac1{(rm)!}\sum_{\pi\in\mathfrak S_{rm}}A_k^\pi(s),
\end{equation}
where $A_k^\pi$ has its bulk values placed according to $\pi$.
Changing variables separately for each permutation in
\eqref{eq:src49} yields
\[
 \eta_k=\E\sum_{0<|s|\le V_k}G_R(s/d)A_{\mathrm{sym}}(s).
\]
In this operation all bins and history-dependent support conditions
remain inside their respective $A_k^\pi$; they need not be invariant.
Cauchy--Schwarz and \eqref{eq:conclusion-regular-norm} give
\begin{equation}\label{eq:conclusion-cauchy-schwarz}
 |\eta_k|^2
 \le\left(\E\sum_s|G_R(s/d)|^2\right)
        \left(\E\sum_s|A_{\mathrm{sym}}(s)|^2\right)
 \ll V_k\E\sum_s|A_{\mathrm{sym}}(s)|^2.
\end{equation}

Expand the second factor into ordered pairs of permutations.  Each
bad pair is at most
$\exp(r(\Delta_0+Cm)+o_k(m))$ in absolute value by Cauchy--Schwarz and
the two one-assignment bounds \eqref{eq:src55}.  Its fraction among all
pairs is bounded by \eqref{eq:conclusion-bad-fraction}.  Each remaining
pair satisfies \eqref{eq:src56}; that estimate is uniform, so averaging
these pairs preserves its bound.  Consequently
\begin{equation}\label{eq:conclusion-final-upper}
 |\eta_k|^2\ll e^{E_k}
 \left[
 \exp\bigl(r\{\Delta_0+(-\tfrac34\log r+C)m\}+o_k(m)\bigr)
       +\exp(-\omega_k(m))
 \right].
\end{equation}

For completeness, the sums in \eqref{eq:src43} satisfy
\[
 \sum_{j=1}^k r_j=r-1,
 \qquad
 \sum_{j=1}^k r_j\log r_j
   =r\log r-2r\log2+2\log2.
\]
They give, in particular, the convenient upper bound
\begin{equation}\label{eq:conclusion-frequency-budget}
 \frac{E_k}{rm}
 \le\frac25\log r+B_D+B_s+(B_z+8)\log z+o_k(1).
\end{equation}
Since $\Delta_0/m=B_s+8\log z$, the logarithm of the first term on the
right of \eqref{eq:conclusion-final-upper}, divided by $rm$, is at most
\begin{equation}\label{eq:conclusion-final-rate}
 -\frac7{20}\log r+B_D+2B_s+(B_z+16)\log z+C+o_k(1).
\end{equation}
All constants here have already been fixed.  Now choose $k$ sufficiently
large, retaining \eqref{eq:conclusion-cg-reserve}, that the expression
in \eqref{eq:conclusion-final-rate} without its $o_k(1)$ term is less
than $-2B-4$.  This is possible because $\log r=k\log2$ while
$\log z=4\log k$.  Finally take $L$ sufficiently large.  The first
term of \eqref{eq:conclusion-final-upper} is then at most
$\exp(-(2B+3)rm)$, after absorbing its absolute implied constant.
The second term remains $\exp(-\omega_k(m))$, since $E_k=O_k(m)$.
Thus, for sufficiently large $L$,
\[
 |\eta_k|^2\le\exp(-(2B+1)rm).
\]
This contradicts \eqref{eq:conclusion-exact-recurrence}, which gives
$|\eta_k|^2\ge\exp(-2Brm)$.  The assumed decomposition is therefore
impossible: no two infinite subsets of $\N_0$ have a sumset whose
symmetric difference with the positive primes is finite.  This proves
Theorem~\ref{thm:main}.
